\documentclass{article}

\usepackage{neurips_2024}

\usepackage[utf8]{inputenc}
\usepackage[T1]{fontenc}
\usepackage{hyperref}
\usepackage{url}
\usepackage{booktabs}
\usepackage{amsfonts}
\usepackage{amsmath}
\usepackage{amssymb}
\usepackage{nicefrac}
%\usepackage{microtype}
\usepackage{xcolor}
\usepackage{algorithm}
\usepackage{algorithmic}
\usepackage{bm}


\title{Beyond Disjoint Tasks: Leveraging Novelty-Learnability Replay and Intrinsic Rewards for Open-Ended Continual Learning}

\author{%
  David S.~Hippocampus \\
  Department of Computer Science\\
  Cranberry-Lemon University\\
  Pittsburgh, PA 15213 \\
  \texttt{hippo@cs.cranberry-lemon.edu} \\
}


\begin{document}

\maketitle

\begin{abstract}
\textcolor{orange}{
\begin{enumerate}
    \item Context: World models (RSSM) provide a powerful framework for RL by training policies in imagination.
    \item The Problem: Current Continual RL (CRL) methods assume "disjoint tasks," but open-ended environments like Craftax feature hierarchical and overlapping task structures where early skills are prerequisites for later ones.
    \item Proposed Innovation: We introduce three enhancements to DreamerV3: (i) NLR/NLU sampling selectors for automated curriculum building; (ii) symbolic intrinsic rewards (Spatial + Craft novelty); and (iii) policy refinements including action masking .
    \item Results: Our approach mitigates the collapse of exploration and enables agents to master complex tech-tree achievements in sparse-reward, open-ended settings.
\end{enumerate}
}

\end{abstract}

%=================================================================
\section{Introduction}
\label{sec:intro}
%=================================================================
\textcolor{orange}{
\begin{enumerate}
    \item Evolution of CRL: Moving from solving isolated tasks sequentially to mastering open-ended domains.
    \item The Challenge of Overlap: In Craftax, state spaces are not disjoint; mastering a "base step" like placing table is essential for maximizing future rewards.
    \item Failure of Baselines: Standard sampling (50:50) and exploration (Plan2Explore) fail because they over-represent early-stage data, leading to over-confidence and exploration collapse.
    \item Contribution Summary: List the three families of innovations (NLR/NLU, Intrinsic Rewards, Action Masking) and their integration into the RSSM.
\end{enumerate}
}


%=================================================================
\section{Related Work}
\label{sec:related_work}
%=================================================================

\textcolor{orange}{
\begin{enumerate}
    \item World Models in RL: Overview of DreamerV2/V3 and RSSM architecture.
    \item Continual RL: Contrast task-aware (EWC) vs. task-agnostic (CLEAR) methods .
    \item Experience Replay: Prioritize the history of Selective Experience Replay (PER, Reservoir Sampling).
    \item Intrinsic Motivation: Discuss P2E and curiosity-driven methods and their limits in open-ended worlds.
\end{enumerate}
}
\subsection{DreamerV3 and the RSSM Architecture}
World models have become an important direction in reinforcement learning (RL) because they enable agents to learn a compact latent representation of the environment and use this learned dynamics model for policy optimization. Model-based RL methods can simulate future trajectories in latent space, improving data efficiency and long-horizon reasoning. A representative line of work is the Dreamer family, which performs behavior learning through imagination rollouts generated by a learned world model \cite{hafner2019dream,hafner2023mastering}.

A key component underlying Dreamer is the Recurrent State-Space Model (RSSM), originally developed in PlaNet \cite{hafner2019learning}. RSSM combines deterministic recurrent hidden states with stochastic latent variables to model partially observable environments. This design allows the model to capture both temporal dependencies and uncertainty in observations. Given an action and the previous latent state, the RSSM predicts a prior distribution over the next latent variable, while an encoder incorporates the current observation to form a posterior state during training. The latent dynamics can then be decoded to reconstruct observations, rewards, and termination signals. By separating deterministic memory from stochastic uncertainty, RSSM provides a structured latent representation that is well suited for sequential decision-making problems \cite{hafner2019learning}.

Building on this architecture, Dreamer introduced actor-critic learning entirely in latent imagination space \cite{hafner2019dream}. After learning the world model from experience, the agent rolls out imagined trajectories inside the RSSM and updates the policy and value networks using these latent predictions, rather than requiring expensive interactions with the real environment for every gradient step. This substantially improves sample efficiency and makes the approach more scalable to visually rich tasks. Dreamer therefore demonstrated that a compact latent world model can support effective control directly from high-dimensional observations \cite{hafner2019dream}.

DreamerV3 further generalizes this framework and shows that a unified world-model-based agent can perform robustly across highly diverse domains with minimal domain-specific tuning \cite{hafner2023mastering}. It preserves the core Dreamer principle of learning behaviors through latent imagination, while improving optimization stability, robustness, and scalability. As a result, DreamerV3 highlights the broader potential of world models as a general-purpose backbone for RL across different observation modalities and task structures \cite{hafner2023mastering}. In this sense, the Dreamer line establishes RSSM-based latent dynamics modeling as a strong foundation for modern model-based RL, and it provides the architectural basis for many subsequent extensions in continual learning, intrinsic motivation, and structured exploration.


%=================================================================
\section{Problem Analysis}
\label{sec:analysis}
%=================================================================
\textcolor{orange}{
\begin{enumerate}
    \item Hierarchical State Overlap: Mathematically define why the $\mathcal{S}_i \cap \mathcal{S}_j = \emptyset$ assumption fails in Craftax.
    \item Sampling Bias: Show how uniform or recency-biased (50:50) sampling wastes computation on mastered "easy" states.
    \item Exploration Collapse: Detail the "P2E collapse" where ensemble variance disappears because the model becomes over-confident in frequently visited early-stage states.
\end{enumerate}
}




%=================================================================
\section{Methodology}
\label{sec:method}
%=================================================================

We extend DreamerV3 with three families of innovations:
(i)~a suite of experience-replay sampling strategies culminating in
\emph{Novelty--Learnability--Recency/Uniform} (NLR/NLU) selectors
that construct an autocurriculum over replay experience
(\S\ref{sec:replay});
(ii)~observation-grounded intrinsic rewards combining spatial-counting
and craft-novelty bonuses that exploit the structured symbolic
observation (\S\ref{sec:intrinsic});
and (iii)~a modified policy that fuses intrinsic and extrinsic
signals during world-model imagination and optionally masks infeasible
actions (\S\ref{sec:policy}).
All components are integrated into DreamerV3's RSSM world model and
trained end-to-end.  Symbolic-observation embedding is described in
\S\ref{sec:embedding}.

%------------------------------------------------------------------
\subsection{Replay Buffer Strategies}
\label{sec:replay}
%------------------------------------------------------------------

Let $\mathcal{B}=\{\tau_1,\dots,\tau_N\}$ be the replay buffer with
capacity~$C$.  Each training step draws a mini-batch of~$B$
trajectories.  We describe first the \emph{eviction} policy (which
items are retained when the buffer is full), then the
\emph{sampling selectors} (which items enter each mini-batch).

Let $\mathcal{B}=\{\tau_1,\dots,\tau_N\}$ be the replay buffer with
capacity~$C$. Each training step draws a mini-batch of~$B$
trajectories. Replay design is particularly important in continual RL,
where retaining diverse historical experience can help mitigate
catastrophic forgetting and improve adaptation across changing tasks
\cite{rolnick2019experience,yang2024augmenting}. We describe first the
\emph{eviction} policy (which items are retained when the buffer is
full), then the \emph{sampling selectors} (which items enter each
mini-batch).

%---- Eviction ----
\subsubsection{Eviction: FIFO vs.\ Reservoir Sampling}
\label{sec:eviction}

\textbf{FIFO (default DreamerV3).}\;
The oldest trajectory is evicted when the buffer reaches capacity.
This is efficient but systematically erases early experience that may
contain first-discovery demonstrations critical for long-horizon skill
chains.

\textbf{Reservoir Sampling.}\;
Following Vitter's Algorithm~R, the $n$-th trajectory arriving at a
full buffer is accepted with probability $C/n$; if accepted it
replaces a uniformly random existing item:
\begin{equation}
  P(\tau_k \in \mathcal{B})
  \;=\;
  \frac{C}{n_{\mathrm{seen}}},
  \qquad \forall\; k \le n_{\mathrm{seen}}.
  \label{eq:reservoir}
\end{equation}
Every past trajectory has equal retention probability regardless of
age.  We use a dedicated RNG and $O(1)$ swap-and-pop so that
eviction adds negligible overhead.

%---- Selectors ----
\subsubsection{Sampling Selectors}
\label{sec:selectors}

Each selector exposes a single callable returning one item key; the
replay buffer invokes it $B$ times per mini-batch.

\paragraph{(a) Uniform.}
Each trajectory is sampled with probability $1/|\mathcal{B}|$.

\paragraph{(b) Recency-weighted.}
A pre-computed triangular distribution assigns weight
$w_i \propto 1 - i/|\mathcal{B}|$ ($i\!=\!0$ is the most recent).
The CDF is stored in a hierarchical tree (branching factor 16) for
$O(\log|\mathcal{B}|)$ sampling.

\paragraph{(c) 50:50 Recent/Uniform Mixture.}
Half of each mini-batch is drawn uniformly; the other half from the
$W$ most recent items with triangular weights.  Implemented via a
general \texttt{Mixture} selector accepting arbitrary sub-selectors
and fractional weights:
\begin{equation}
  P(\tau_k)
  = \tfrac12\;\text{Uniform}(\tau_k)
  + \tfrac12\;\text{Recency}_W(\tau_k).
  \label{eq:fiftyfifty}
\end{equation}

\paragraph{(d) Reward-weighted.}
Trajectories are sampled proportionally to a softmax over cumulative
returns:
\begin{equation}
  P(\tau_k)
  =
  \frac{\exp(R_k/T)}
       {\sum_{j}\exp(R_j/T)},
  \qquad
  R_k = {\textstyle\sum_t} r_t^{(k)},
  \label{eq:reward_sel}
\end{equation}
where $T$ is a temperature.  Returns are updated lazily once each
episode completes.

\paragraph{(e) Novelty--Learnability--Recency (NLR) \& NLU.}
\emph{This is our primary novel contribution.}
Rather than treating the buffer as homogeneous, we partition each
mini-batch into three \emph{pools}:
\begin{enumerate}
\item \textbf{Novelty pool} ($\beta_n B$ samples)~---
  oversamples trajectories containing rarely accomplished skills.
\item \textbf{Learnability pool} ($\beta_l B$ samples)~---
  prioritises trajectories whose return exceeds a running baseline
  (``hard but solvable'').
\item \textbf{Third pool} ($\beta_r B$ samples)~---
  \textbf{NLR}: triangular recency over a window of size~$W$.\;
  \textbf{NLU}: uniform over the entire buffer.
\end{enumerate}
Pool fractions satisfy $\beta_n+\beta_l+\beta_r=1$
(defaults $0.35,\,0.35,\,0.30$).
A single uniform draw selects the pool; a weighted sample is then
drawn within it (Algorithm~\ref{alg:nlr}).

\begin{algorithm}[t]
\caption{NLR/NLU Replay Selector --- single item draw}
\label{alg:nlr}
\begin{algorithmic}[1]
\REQUIRE Buffer keys $\mathcal{K}$, pool fractions
  $(\beta_n,\beta_l,\beta_r)$, temperatures $(\tau_n,\tau_l)$,
  window $W$
\STATE $r \sim \mathrm{Uniform}(0,1)$
\IF{$r < \beta_n$ \AND novelty pool $\neq\emptyset$}
  \STATE \textbf{return} key $\sim$
    $s_{\mathrm{novel}}(k)^{1/\tau_n}$
\ELSIF{$r < \beta_n{+}\beta_l$ \AND learnability pool $\neq\emptyset$}
  \STATE \textbf{return} key $\sim$
    $s_{\mathrm{learn}}(k)^{1/\tau_l}$
\ELSE
  \STATE \textbf{return} key from third pool
    {\small(NLR: triangular over last $W$;\; NLU: uniform)}
\ENDIF
\end{algorithmic}
\end{algorithm}

We provide two scoring families---\emph{privileged} and
\emph{non-privileged}---that share the same learnability scorer and
pool-selection logic but differ in how novelty is computed.

\medskip\noindent
\textbf{Privileged NLR/NLU.}\quad
Let the environment expose a binary achievement vector
$\bm{a}^{(k)}\!\in\!\{0,1\}^A$ per episode (Craftax: $A\!=\!67$).
We maintain a running per-achievement success rate
$p_a = (\text{\# episodes achieving }a) / (\text{\# total episodes})$.
The \emph{novelty score} for $\tau_k$ with achieved set
$\mathcal{A}_k=\{a:a_a^{(k)}\!=\!1\}$ is:
\begin{equation}
  s_{\mathrm{novel}}(\tau_k)
  =
  \frac{1}{|\mathcal{A}_k|}
  \sum_{a\in\mathcal{A}_k}
  \frac{1}{p_a+\epsilon}\,.
  \label{eq:novelty_priv}
\end{equation}
Rarely accomplished achievements produce large $1/(p_a{+}\epsilon)$
terms, steering the novelty pool toward under-practised skills.
Trajectories with $\mathcal{A}_k\!=\!\emptyset$ score zero and never
enter the pool.

The \emph{learnability score} implements GRPO-style advantage
filtering:
\begin{equation}
  s_{\mathrm{learn}}(\tau_k)
  =
  \max\!\bigl(0,\; R_k - \bar{R}\bigr),
  \label{eq:learnability}
\end{equation}
where $\bar{R}$ is an EMA of episodic returns with decay
$\gamma_{\mathrm{ema}}$.  Only episodes with positive advantage
enter the pool.

\emph{Note:} the privileged variant uses per-achievement annotations
that the policy network \emph{never observes}; the annotations only
influence \emph{which} trajectories are sampled.

\medskip\noindent
\textbf{Non-privileged NLR/NLU.}\quad
When per-achievement information is unavailable we estimate novelty
from a quantile-adaptive 2-D histogram over
(episode length~$L$, episodic reward~$R$).
The space is partitioned into $B_L\times B_R$ bins with
quantile-based edges (recomputed every $N_{\mathrm{recomp}}$
episodes).  For $\tau_k$ in bin $(b_L,b_R)$ with midpoint reward
$\hat{R}_{b_R}$ and bin count $n_{b_L,b_R}$:
\begin{equation}
  s_{\mathrm{novel}}(\tau_k)
  =
  \frac{
    \sigma\!\!\left(\frac{\hat{R}_{b_R}-R_{\min}}{\beta}\right)
    \cdot\hat{R}_{b_R}
  }{
    n_{b_L,b_R}+\epsilon
  }\,,
  \label{eq:novelty_nonpriv}
\end{equation}
where $R_{\min}$ is the $q$-th percentile of observed rewards
(default $q\!=\!0.2$),
$\beta=\max(0.1,(R_{50}-R_{\min})/4)$ is an adaptive sharpness,
and $\sigma$ is the logistic sigmoid.
The $1/(n+\epsilon)$ factor promotes \emph{rarity} while the
$\sigma(\cdot)\cdot\hat{R}$ prior suppresses low-reward bins,
preventing the novelty pool from being flooded by plentiful but
uninformative short/low-reward episodes.
Learnability scoring and pool selection are identical to the
privileged variant.  The pool-update procedure executed at each
episode boundary is given in Algorithm~\ref{alg:pool_update}
(Appendix~\ref{app:pool_update}).


%------------------------------------------------------------------
\subsection{Intrinsic Reward Design}
\label{sec:intrinsic}
%------------------------------------------------------------------

Open-ended environments have extremely sparse extrinsic reward.  We
introduce two complementary observation-grounded intrinsic bonuses
that exploit the structure of the symbolic observation
$o_t=[v_{\mathrm{vis}}^{(t)};\,v_{\mathrm{inv}}^{(t)}]$
(local map $9{\times}11{\times}83$ plus 51-element inventory).

\paragraph{Spatial-counting reward.}
Define a hash $\phi:\mathbb{R}^{d_{\mathrm{vis}}}\!\to\!\mathbb{Z}^K$
discretising the visual state and maintain an episodic count table
$N(\phi(s))$:
\begin{equation}
  r_{\mathrm{spatial}}^{(t)}
  =
  \frac{1}{\sqrt{N\bigl(
    \phi(v_{\mathrm{vis}}^{(t+1)})\bigr)}}\,.
  \label{eq:spatial}
\end{equation}
As the agent loops, $N\!\to\!\infty$ and $r\!\to\!0$, providing
diminishing incentive to revisit known locations.

\paragraph{Craft-novelty reward.}
Reward only the \emph{first discovery} of a specific inventory
configuration.  Let $\mathcal{H}_{\mathrm{inv}}$ be the set of all
previously seen inventory vectors:
\begin{equation}
  r_{\mathrm{craft}}^{(t)}
  =
  \mathbb{I}\!\bigl[
    v_{\mathrm{inv}}^{(t+1)}\notin\mathcal{H}_{\mathrm{inv}}
  \bigr].
  \label{eq:craft}
\end{equation}
This sparse, non-decaying signal forces the agent to advance along
the tech tree to receive further reward.

\paragraph{Composite intrinsic reward.}
The two components are combined:
\begin{equation}
\boxed{\;
  r_t^{\mathrm{intr}}
  =
  r_{\mathrm{spatial}}^{(t)}
  +
  \lambda\,r_{\mathrm{craft}}^{(t)}
\;}
  \label{eq:intr_total}
\end{equation}
where $\lambda$ weights craft-novelty relative to spatial counting
(default 1.0).  The spatial term discourages within-episode looping
by penalising revisits to known map cells; the craft term provides
directed incentive to climb the technology tree by rewarding
first-discovery of novel inventory configurations.


%------------------------------------------------------------------
\subsection{Policy Architecture}
\label{sec:policy}
%------------------------------------------------------------------

\subsubsection{Combined Reward in Imagination}
\label{sec:combined_reward}

DreamerV3's actor-critic trains entirely inside the world model.
At each imagined step the reward head predicts extrinsic return
$r_t^{\mathrm{extr}}$; we combine it with the intrinsic signal:
\begin{equation}
  r_t
  =
  \alpha_i\,r_t^{\mathrm{intr}}
  +
  \alpha_e\,r_t^{\mathrm{extr}},
  \label{eq:combined}
\end{equation}
with defaults $\alpha_i\!=\!\alpha_e\!=\!0.9$.  This blended reward
enters DreamerV3's standard $\lambda$-return with percentile
normalisation.  The actor-critic thus learns to seek both novel states
(via intrinsic bonuses) and environment progress (via extrinsic
reward) simultaneously.

\subsubsection{Action Masking}
\label{sec:masking}

Craftax exposes 43 discrete actions, but many are infeasible in a
given state (e.g.\ \textsc{place\_table} requires $\ge 2$ wood and a
placeable block ahead).  Rather than hard-blocking actions, we
introduce a \emph{soft feasibility bias} on policy logits that
smoothly discourages infeasible actions while preserving gradient flow.

\paragraph{Context extraction.}
At each environment step we extract a compact
\emph{action-mask context} vector
$c_t\!\in\!\mathbb{R}^{46}$ from the game state.  Its entries encode
inventory counts (wood, stone, coal, iron, diamond, \dots), player
attributes (health, energy, mana), proximity flags (adjacent to
crafting table or furnace), facing-block type, and ladder/level
indicators.  This extraction runs in JAX and is batched across all
parallel environments.

\paragraph{Declarative feasibility rules.}
We define a rule set $\mathcal{R}=\{(a,\,\mathcal{C}_a)\}$ mapping
37 maskable actions to conjunctions of typed conditions over~$c_t$.
Four primitive condition types suffice:
\begin{itemize}
  \item \texttt{min}$(i,v)$: requires $c_t[i]\!\ge\!v$\;
        (e.g.\ wood $\ge 2$),
  \item \texttt{bool}$(i)$: requires $c_t[i]\!>\!0.5$\;
        (e.g.\ near crafting table),
  \item \texttt{below}$(i,v)$: requires $c_t[i]\!<\!v$,
  \item \texttt{any\_below}$(i{:}j,\,v)$: requires at least one slot
        in range $[i,j)$ to satisfy $<v$.
\end{itemize}
For each action~$a$ we compute a \emph{deficit} that aggregates how
far the current context is from satisfying all conditions:
\begin{equation}
  d_a(c_t)
  =
  \sum_{(\texttt{min},\,i,\,v)\,\in\,\mathcal{C}_a}
    \max\!\bigl(0,\; v - c_t[i]\bigr)
  \;+\;
  \sum_{(\texttt{bool},\,i)\,\in\,\mathcal{C}_a}
    \mathbb{I}\!\bigl[c_t[i]\le 0.5\bigr].
  \label{eq:deficit}
\end{equation}
When all conditions are satisfied, $d_a\!=\!0$; each unmet condition
contributes a positive term proportional to its violation magnitude.

\paragraph{Soft-mode logit bias.}
The deficit is converted to a logit bias that is \emph{added} to the
raw policy logits $\ell_a$ before the softmax:
\begin{equation}
  b_a
  =
  -\lambda\,d_a(c_t),
  \qquad
  \pi_{\mathrm{mask}}(a\mid s)
  =
  \frac{\exp(\ell_a + b_a)}
       {\sum_{a'}\exp(\ell_{a'} + b_{a'})},
  \label{eq:softmask}
\end{equation}
where $\lambda$ is a penalty strength (default $\lambda\!=\!5$).
Unlike binary masking, the soft bias (i)~preserves non-zero
probability on every action, avoiding gradient death for infeasible
actions; (ii)~provides a graded signal---actions that are ``almost
feasible'' (low deficit) receive a milder penalty than clearly
impossible ones; and (iii)~is differentiable with respect to the
context, which matters when the context is predicted rather than
observed (see below).

For comparison we also implement a \textbf{hard mode} that sets
$b_a = -10^9$ whenever $d_a>0$, effectively zeroing the probability
of any infeasible action.  Empirically, hard masking destabilises
training (Section~\ref{sec:experiments}).

\paragraph{Mask context prediction for imagination.}
During real environment interaction, the true context $c_t$ is
available.  However, DreamerV3's actor-critic trains inside
\emph{imagined} rollouts where no true game state exists.  To enable
masking during imagination we train an auxiliary
\emph{mask-context head} $f_\psi$ that predicts $c_t$ from RSSM
features:
\begin{equation}
  \hat{c}_t = f_\psi\!\bigl(\mathrm{sg}(z_t)\bigr),
  \qquad
  \mathcal{L}_{\mathrm{ctx}}
  =
  -\!\sum_{t}\log p_\psi(c_t \mid z_t),
  \label{eq:ctx_head}
\end{equation}
where $z_t$ denotes the concatenated deterministic and stochastic RSSM
state, $\mathrm{sg}(\cdot)$ is the stop-gradient operator, and
$p_\psi$ is a diagonal Gaussian with learned mean and log-standard
deviation (a 2-layer MLP with SiLU activations and RMSNorm).
The head is trained jointly with the world model from step~0,
with loss scale 1.0.

\paragraph{Dream-mask warmup.}
Early in training, $f_\psi$ predictions are unreliable.  We therefore
activate masking in imagination only after a warmup period of $W_m$
steps (default $W_m\!=\!50\text{k}$, including 10k prefill).
Before this point, imagination rollouts use the unmasked policy;
afterwards, each imagined step computes
$\hat{c}_t=f_\psi(z_t)$, applies the soft bias from
Eq.~\eqref{eq:softmask}, and samples from the resulting adjusted
distribution.  The warmup is implemented via a Python-level flag that
triggers JAX re-compilation when flipped.

\subsubsection{Separate Exploration and Exploitation Policies}
\label{sec:sep_policies}

Optionally we maintain two policy heads sharing the world model:
an \emph{exploration policy} $\pi^{\mathrm{expl}}$ trained with
the combined reward
$r_t=\alpha_i r_t^{\mathrm{intr}}+\alpha_e r_t^{\mathrm{extr}}$,
and an \emph{evaluation policy} $\pi^{\mathrm{eval}}$ trained with
$r_t=r_t^{\mathrm{extr}}$ only.  The exploration policy collects
data; the evaluation policy is used for assessment.  Both share the
RSSM, differing only in their MLP actor-critic heads, thereby
decoupling the exploration objective from the evaluation metric.


%------------------------------------------------------------------
\subsection{Embedding Mode for Symbolic Observations}
\label{sec:embedding}
%------------------------------------------------------------------

Craftax's symbolic observation is a flat 8\,268-dimensional vector
(a $9{\times}11{\times}83$ map tensor plus 51-element inventory).
Rather than using DreamerV3's convolutional encoder (designed for
pixels), we project into a compact embedding:
\begin{equation}
  e_t
  =
  W_{\mathrm{enc}}\,\mathrm{flatten}(o_t)+b_{\mathrm{enc}},
  \qquad
  e_t\in\mathbb{R}^d,
  \label{eq:embedding}
\end{equation}
with $d\!=\!256$.  This replaces the convolutional encoder output and
feeds directly to the RSSM; the decoder is similarly replaced with an
MLP.  Eliminating image reconstruction lets the RSSM focus on dynamics
and reward prediction, yielding faster, more stable training.





%=================================================================
\section{Experiments}
\label{sec:experiments}
%=================================================================

\textcolor{orange}{
\begin{enumerate}
    \item Setup: Describe Craftax-symbolic-v1 with its 67-dimensional achievement vector.
    \item Comparative Performance: Compare DreamerV3-CL against Original DreamerV3, highlighting the achievement rates of critical items like placing table.
    \item Loss Analysis: Interpret why Dynamic and Representation losses spike in the CL version (indicating higher confidence in complex transitions).
    \item Ablation Studies: Test the contribution of the Novelty vs. Learnability pools and the impact of Action Masking on training stability.
\end{enumerate}
}

%=================================================================
\section{Conclusion and Future Work}
\label{sec:conclusion}
%=================================================================





%%%%%%%%%%%%%%%%%%%%%%%%%%%%%%%%%%%%%%%%%%%%%%%%%%%%%%%%%%%%

\appendix

\section{Supplementary Methodology Details}

\subsection{NLR Pool Update Procedure}
\label{app:pool_update}

Algorithm~\ref{alg:pool_update} describes how the novelty and
learnability pools are updated at each episode boundary.

\begin{algorithm}[h]
\caption{NLR pool update (per completed episode)}
\label{alg:pool_update}
\begin{algorithmic}[1]
\REQUIRE Episode return $R_k$, achievements $\bm{a}^{(k)}$
  (privileged) or length $L_k$ (non-priv.),
  item keys $\{k_j\}$ inserted during episode
\STATE $\bar{R}\leftarrow\gamma_{\mathrm{ema}}\bar{R}
  +(1{-}\gamma_{\mathrm{ema}})R_k$
\STATE \textit{// Novelty pool}
\IF{privileged}
  \STATE Update $p_a$ for all $a$;\;
    $s_{\mathrm{novel}}\leftarrow$Eq.~\eqref{eq:novelty_priv}
\ELSE
  \STATE Assign $(L_k,R_k)$ to bin $(b_L,b_R)$;\;
    $n_{b_L,b_R}\mathrel{{+}{=}}1$;\;
    $s_{\mathrm{novel}}\leftarrow$Eq.~\eqref{eq:novelty_nonpriv}
  \STATE Every $N_{\mathrm{recomp}}$ eps: recompute quantile edges;\;
    rebuild all scores
\ENDIF
\IF{$s_{\mathrm{novel}}>0$}
  \STATE Add $\{k_j\}$ to novelty pool with $s_{\mathrm{novel}}$
\ENDIF
\STATE \textit{// Learnability pool}
\STATE $s_{\mathrm{learn}}\leftarrow\max(0,R_k{-}\bar{R})$
\IF{$s_{\mathrm{learn}}>0$}
  \STATE Add $\{k_j\}$ to learnability pool with $s_{\mathrm{learn}}$
\ENDIF
\end{algorithmic}
\end{algorithm}

\subsection{Full Hyperparameter Table}
\label{app:hyperparams}

\begin{table}[h]
\centering
\caption{Default hyperparameters for all proposed components.}
\label{tab:hyperparams}
\small
\begin{tabular}{@{}p{3.4cm}p{4.8cm}p{3.0cm}@{}}
\toprule
\textbf{Component} & \textbf{Parameter} & \textbf{Default} \\
\midrule
\multicolumn{3}{@{}l}{\textit{Replay Buffer}} \\
& Capacity $C$ & $2\!\times\!10^6$ \\
& Eviction & Reservoir \\
\midrule
\multicolumn{3}{@{}l}{\textit{NLR / NLU Pools}} \\
& $\beta_n,\;\beta_l,\;\beta_r$ & 0.35,\;0.35,\;0.30 \\
& Recency window $W$ & 1000 \\
& EMA decay $\gamma_{\mathrm{ema}}$ & 0.99 \\
& Temperatures $\tau_n,\tau_l$ & 1.0,\;1.0 \\
\midrule
\multicolumn{3}{@{}l}{\textit{Privileged Only}} \\
& Smoothing $\epsilon$ & 0.01 \\
\midrule
\multicolumn{3}{@{}l}{\textit{Non-Priv.\ Grid}} \\
& Bins $B_R\!\times\!B_L$ & $5\!\times\!10$ \\
& Recomp.\ interval & 500 eps \\
& Prior percentile $q$ & 0.20 \\
& Grid $\epsilon$ & 0.01 \\
\midrule
\multicolumn{3}{@{}l}{\textit{Combined Reward}} \\
& $\alpha_i,\;\alpha_e$ & 0.9,\;0.9 \\
& Craft weight $\lambda$ & 1.0 \\
\midrule
\multicolumn{3}{@{}l}{\textit{Embedding}} \\
& Dimension $d$ & 256 \\
\midrule
\multicolumn{3}{@{}l}{\textit{DreamerV3 (size25m)}} \\
& RSSM deter / hidden & 3072 / 384 \\
& stoch $\times$ classes & 32 $\times$ 24 \\
& Imagination horizon & 15 \\
& Batch / seq.\ length & 16 / 64 \\
& Parallel envs & 64 \\
& Optimiser & Adam, lr $4{\times}10^{-5}$ \\
\bottomrule
\end{tabular}
\end{table}

\subsection{Non-Privileged Grid Recomputation}
\label{app:grid}

The quantile-adaptive 2-D grid is rebuilt every $N_{\mathrm{recomp}}$
episodes:
(1)~Collect all $(L_k,R_k)$ pairs in the buffer.
(2)~Compute $B_L{+}1$ and $B_R{+}1$ quantile edges for the length
  and reward marginals; perturb coinciding edges by $10^{-6}$.
(3)~Set midpoints
  $\hat{R}_b=(\text{edge}_b{+}\text{edge}_{b+1})/2$.
(4)~Set $R_{\min}$ to the $q$-th percentile;
  $\beta=\max(0.1,(R_{50}{-}R_{\min})/4)$.
(5)~Re-assign all trajectories to new bins; rebuild counts and
  novelty scores.
Between recomputations, new episodes use existing edges with
incremental count updates.


%%%%%%%%%%%%%%%%%%%%%%%%%%%%%%%%%%%%%%%%%%%%%%%%%%%%%%%%%%%%

\newpage
\section*{NeurIPS Paper Checklist}

%%% BEGIN INSTRUCTIONS %%%
The checklist is designed to encourage best practices for responsible machine learning research, addressing issues of reproducibility, transparency, research ethics, and societal impact. Do not remove the checklist: {\bf The papers not including the checklist will be desk rejected.} The checklist should follow the references and follow the (optional) supplemental material.  The checklist does NOT count towards the page
limit. 

Please read the checklist guidelines carefully for information on how to answer these questions. For each question in the checklist:
\begin{itemize}
    \item You should answer \answerYes{}, \answerNo{}, or \answerNA{}.
    \item \answerNA{} means either that the question is Not Applicable for that particular paper or the relevant information is Not Available.
    \item Please provide a short (1–2 sentence) justification right after your answer (even for NA). 
\end{itemize}

{\bf The checklist answers are an integral part of your paper submission.} They are visible to the reviewers, area chairs, senior area chairs, and ethics reviewers. You will be asked to also include it (after eventual revisions) with the final version of your paper, and its final version will be published with the paper.

The reviewers of your paper will be asked to use the checklist as one of the factors in their evaluation. While "\answerYes{}" is generally preferable to "\answerNo{}", it is perfectly acceptable to answer "\answerNo{}" provided a proper justification is given (e.g., "error bars are not reported because it would be too computationally expensive" or "we were unable to find the license for the dataset we used"). In general, answering "\answerNo{}" or "\answerNA{}" is not grounds for rejection. While the questions are phrased in a binary way, we acknowledge that the true answer is often more nuanced, so please just use your best judgment and write a justification to elaborate. All supporting evidence can appear either in the main paper or the supplemental material, provided in appendix. If you answer \answerYes{} to a question, in the justification please point to the section(s) where related material for the question can be found.

IMPORTANT, please:
\begin{itemize}
    \item {\bf Delete this instruction block, but keep the section heading ``NeurIPS paper checklist"},
    \item  {\bf Keep the checklist subsection headings, questions/answers and guidelines below.}
    \item {\bf Do not modify the questions and only use the provided macros for your answers}.
\end{itemize} 

%%% END INSTRUCTIONS %%%


\begin{enumerate}

\item {\bf Claims}
    \item[] Question: Do the main claims made in the abstract and introduction accurately reflect the paper's contributions and scope?
    \item[] Answer: \answerTODO{}
    \item[] Justification: \justificationTODO{}

\item {\bf Limitations}
    \item[] Question: Does the paper discuss the limitations of the work performed by the authors?
    \item[] Answer: \answerTODO{}
    \item[] Justification: \justificationTODO{}

\item {\bf Theory Assumptions and Proofs}
    \item[] Question: For each theoretical result, does the paper provide the full set of assumptions and a complete (and correct) proof?
    \item[] Answer: \answerTODO{}
    \item[] Justification: \justificationTODO{}

\item {\bf Experimental Result Reproducibility}
    \item[] Question: Does the paper fully disclose all the information needed to reproduce the main experimental results of the paper to the extent that it affects the main claims and/or conclusions of the paper (regardless of whether the code and data are provided or not)?
    \item[] Answer: \answerTODO{}
    \item[] Justification: \justificationTODO{}

\item {\bf Open access to data and code}
    \item[] Question: Does the paper provide open access to the data and code, with sufficient instructions to faithfully reproduce the main experimental results, as described in supplemental material?
    \item[] Answer: \answerTODO{}
    \item[] Justification: \justificationTODO{}

\item {\bf Experimental Setting/Details}
    \item[] Question: Does the paper specify all the training and test details (e.g., data splits, hyperparameters, how they were chosen, type of optimizer, etc.) necessary to understand the results?
    \item[] Answer: \answerTODO{}
    \item[] Justification: \justificationTODO{}

\item {\bf Experiment Statistical Significance}
    \item[] Question: Does the paper report error bars suitably and correctly defined or other appropriate information about the statistical significance of the experiments?
    \item[] Answer: \answerTODO{}
    \item[] Justification: \justificationTODO{}

\item {\bf Experiments Compute Resources}
    \item[] Question: For each experiment, does the paper provide sufficient information on the computer resources (type of compute workers, memory, time of execution) needed to reproduce the experiments?
    \item[] Answer: \answerTODO{}
    \item[] Justification: \justificationTODO{}

\item {\bf Code Of Ethics}
    \item[] Question: Does the research conducted in the paper conform, in every respect, with the NeurIPS Code of Ethics \url{https://neurips.cc/public/EthicsGuidelines}?
    \item[] Answer: \answerTODO{}
    \item[] Justification: \justificationTODO{}

\item {\bf Broader Impacts}
    \item[] Question: Does the paper discuss both potential positive societal impacts and negative societal impacts of the work performed?
    \item[] Answer: \answerTODO{}
    \item[] Justification: \justificationTODO{}

\item {\bf Safeguards}
    \item[] Question: Does the paper describe safeguards that have been put in place for responsible release of data or models that have a high risk for misuse (e.g., pretrained language models, image generators, or scraped datasets)?
    \item[] Answer: \answerTODO{}
    \item[] Justification: \justificationTODO{}

\item {\bf Licenses for existing assets}
    \item[] Question: Are the creators or original owners of assets (e.g., code, data, models), used in the paper, properly credited and are the license and terms of use explicitly mentioned and properly respected?
    \item[] Answer: \answerTODO{}
    \item[] Justification: \justificationTODO{}

\item {\bf New Assets}
    \item[] Question: Are new assets introduced in the paper well documented and is the documentation provided alongside the assets?
    \item[] Answer: \answerTODO{}
    \item[] Justification: \justificationTODO{}

\item {\bf Crowdsourcing and Research with Human Subjects}
    \item[] Question: For crowdsourcing experiments and research with human subjects, does the paper include the full text of instructions given to participants and screenshots, if applicable, as well as details about compensation (if any)?
    \item[] Answer: \answerTODO{}
    \item[] Justification: \justificationTODO{}

\item {\bf Institutional Review Board (IRB) Approvals or Equivalent for Research with Human Subjects}
    \item[] Question: Does the paper describe potential risks incurred by study participants, whether such risks were disclosed to the subjects, and whether Institutional Review Board (IRB) approvals (or an equivalent approval/review based on the requirements of your country or institution) were obtained?
    \item[] Answer: \answerTODO{}
    \item[] Justification: \justificationTODO{}

\end{enumerate}

\bibliographystyle{plainnat}
\bibliography{Reference}


\end{document}